\chapter{REVISÃO DA LITERATURA}

Através da revisão da literatura, busca-se contextualizar o desenvolvimento do sistema proposto para a loja \textbf{Tropical Mix}, destacando a relevância da automação de processos comerciais, do controle de estoque e da digitalização de operações em pequenos e médios empreendimentos. Assim, pretende-se apresentar o embasamento teórico necessário para a compreensão do problema e das soluções tecnológicas associadas.

\section{Histórico da Gestão Automatizada}

A informatização de processos empresariais teve início com a popularização dos sistemas de gestão nos anos 1980, quando grandes corporações passaram a utilizar softwares para controle de produção e finanças, conhecidos como \textit{Enterprise Resource Planning} (ERP). Com o passar dos anos, esses sistemas evoluíram e começaram a atender também pequenas e médias empresas, proporcionando maior controle sobre estoques, vendas e relacionamento com clientes \cite{laudon2022}.
.

Inicialmente, os sistemas eram instalados localmente (\textit{on-premise}), exigindo infraestrutura e manutenção interna. Entretanto, com o avanço da computação em nuvem e a consolidação de serviços baseados na \textit{web}, emergiu uma geração de soluções que combinam acesso remoto — \textit{anytime, anywhere} — com a centralização do armazenamento e da gestão de dados na infraestrutura do provedor \cite[p.~5]{armbrust2009}. Em linha com essa visão, a definição do NIST caracteriza a nuvem como um modelo que habilita acesso ubíquo e sob demanda, via rede, a um conjunto compartilhado de recursos computacionais \cite{mell2011}.
.

No contexto do comércio varejista, a adoção de sistemas de controle de estoque e vendas passou a representar não apenas um diferencial competitivo, mas uma necessidade operacional. Tais ferramentas garantem precisão nas informações, redução de perdas e agilidade na tomada de decisões, fatores essenciais para o crescimento de uma empresa.

\section{GERENCIAMENTO NA ATUALIDADE}




\section{Gerenciamento em Diversos Setores}

